\section{Clifford--Ore表示の外に出る代数的な二周期変形}\label{sec:nonore}
固定された三生成部分代数のOre拡大とは異なる変形を得るため、
三生成元の関係自体に第四生成元$X$を含む項を加える。
ここでは記号$k_0,v_0,b_0,s_0,a_0$をパラメータとし、
基底体$k$との混同を避ける。$S_5=\mathbb A^5_{k_0,v_0,b_0,s_0,a_0}$とする。
位置$i$で$\varepsilon=(-1)^i$と置き、以下の関係を課す。
\begin{equation}\label{eq:nonore-five}
\begin{aligned}
 YW+WY&=\varepsilon a_0ZX,& ZW+WZ&=0,& ZY+YZ&=k_0W^2,\\
 XW+WX&=v_0Z^2,& XY+YX&=b_0W^2+s_0Z^2,& XZ+ZX&=0.
\end{aligned}
\end{equation}
$k_0=v_0=b_0=s_0=1,a_0=0$でClifford基点$C_0$になる。
$a_0$の変化は$YW+WY$に$X$を含ませるので、
\emph{選んだ}$W,Y,Z$部分代数によるOre構成には属さない。

\begin{theorem}[全五変数で正則な非Ore型チャート]\label{thm:nonore-five}
式\eqref{eq:nonore-five}は$S_5$上、全次数で自由な二周期$\ZZ$-代数を定める。
各次数$n$で$W^iY^jZ^\ell X^m$（$i+j+\ell+m=n$）が基底となり、
全ての幾何学的ファイバーは標準14641型のquadratic Koszul AS正則代数である。
添字保持同型による同型類の像の次元は正確に\emph{2}である。
Clifford基点での一次変形は独立な偶方向1個と交代方向1個を含み、
後者は前節の九変数Ore族の接像には入らない。
一般ファイバーの最小周期が2であるかどうかは本定理では主張しない。
\end{theorem}
\begin{proof}
$W<Y<Z<X$とし、補助重みを
\[
 \operatorname{wt}(W,Y,Z,X)=(3,5,2,2)
\]
と定め、正の重みと辞書式順序で語を比較する。
式\eqref{eq:nonore-five}の六つの規則では、反順序の二次語を
$-1$倍の順序付き語へ換える項は同じ重みで辞書式順序が小さく、
残る全ての項は\emph{真に小さい重み}を持つ。
例えば$YW$の重みは8だが$ZX$は4、$ZY$は7だが$W^2$は6、
$XW$は5だが$Z^2$は4、$XY$は7だが$W^2$は6、$Z^2$は4である。
従って還元は全パラメータ上で停止する。

三次重なりは$ZYW$、$XYW$、$XZW$、$XZY$の四つであり、
始点の偶奇を含めて全部で8個である。
$\mathbb Q[k_0,v_0,b_0,s_0,a_0]$上で、二つの還元順序の差は
\emph{恒等的に零}になる。
再現コード\path{computations/verify_nonore_family.py}は8個すべてを
記号多項式として検査する。Diamond lemmaより全基底上の自由性が従う。
重み付き随伴次数代数は、四つの生成元が相異なるとき反可換する
標準四変数量子多項式代数である。
二次Gr\"obner性からKoszul性が従い、二次双対のHilbert級数は
$(1+t)^4$となる。双対への逆フィルトレーションを使うと、
その随伴次数代数は量子外積代数からの商であり、次数別次元が等しいため同型となる。
量子外積の相補次数積は完全であるから元の双対も局所Frobeniusである。
左右Koszul解消を双対化して標準14641型の左右AS条件と大域次元4を得る。
係数は位置の偶奇だけに依存するので厳密2周期的である。

次元を計算する。$S_5$には四生成元の対角変更
$W\mapsto\alpha W,Y\mapsto\beta Y,Z\mapsto\gamma Z,X\mapsto\delta X$
が作用し、係数は
\[
 \begin{aligned}
 k_0&\mapsto k_0\alpha^2/(\beta\gamma),&
 v_0&\mapsto v_0\gamma^2/(\alpha\delta),\\
 b_0&\mapsto b_0\alpha^2/(\beta\delta),&
 s_0&\mapsto s_0\gamma^2/(\beta\delta),&
 a_0&\mapsto a_0\gamma\delta/(\alpha\beta)
 \end{aligned}
\]
と変化する。共通スカラーを除いた有効なトーラス作用の階数は3なので、
その一般軌道は3次元である。
従って同型類の像の次元は高々$5-3=2$である。
このトーラスの有理不変量として、例えば
\[
 I_1=\frac{b_0s_0}{k_0^2v_0^2},\qquad
 I_2=\frac{a_0b_0^4}{k_0^5v_0^3}
\]
を取れる。これらが既に\emph{全ての}indexed同型類を分類するとは主張しない。

逆向きの次元不等式はClifford基点での厳密線型代数から従う。
一つの関係空間上の$\PGL_4\times\PGL_4$軌道微分は階数30、
五パラメータを加えた微分は階数32である。
その偶方向は軌道と合わせて階数16、交代方向も階数16である。
軌道飽和像の次元は少なくとも32で、一般軌道は30次元だから
同型類の像の次元は少なくとも2である。
以上により正確に2となる。

さらに$a_0$の接方向は$YW$の書換え則へ$(-1)^iZX$を加えるだけで、
先に固定した交代商基底$o_0$そのものである。
九変数Ore族の交代接像は
$(o_0,o_1,o_2,o_3,o_4,o_5)=(0,-b,-a,-b,0,c)$なので、
$o_0$はその像に含まれない。
これはClifford基点での\emph{実際の新しい代数的変形方向}を証明する。
ただし他のOre分解を用いた表示や、全モジュライ成分への所属については未決定である。
\end{proof}

\begin{remark}
式\eqref{eq:nonore-five}では一般の$k_0,v_0,b_0,s_0$が零でもAS正則性は成立する。
これは九変数楕円族の$1+th\ne0$という制限とは別の全パラメータ主張である。
一方、二周期的表示を持つことから最小周期2とは結論できない。
一般の最小周期や、九変数族の別のチャートとの大域的な交わりは未解決である。
\end{remark}
